\begin{enumerate}
  \item \textbf{Ruch w VR jest silną i stabilną w czasie cechą biometryczną.} Już modele bazowe
        rozpoznają osobę następnego dnia z Rank-1 0{,}93--0{,}94 (ResNet1D, Transformer), a po
        dopasowaniu treningu do oceny cosinusowej -- 0{,}977 na odłożonym teście (ensemble) przy EER
        ok.~1\%. Kształt rzutu tej samej osoby jest powtarzalny między dniami, a różnice
        międzyosobnicze (tempo, amplituda, orientacja dłoni i głowy) są wyraźne.
  \item \textbf{Reprezentacja danych waży więcej niż wybór architektury.} Dwie największe poprawy w
        projekcie nie wynikają z modelu, tylko z cech wejściowych i funkcji straty:
        (a) ciągła reprezentacja 6D rotacji zamiast kątów Eulera -- odwróciła ranking modalności i
        podniosła Rank-1 rotacji GRU z 0{,}52 do 0{,}74; (b) ArcFace i prędkości -- podniosły
        GRU z 0{,}82 do 0{,}97--0{,}98 i zrównały wszystkie cztery architektury (0{,}960--0{,}972).
        RSA prowadzi do tego samego wniosku: embeddingi różnych modeli z tej samej modalności są do
        siebie bardziej podobne niż embeddingi jednego modelu z różnych modalności.
  \item \textbf{Modalności:} pełny $>$ rotacja $>$ pozycja dla każdego modelu i każdej metryki.
        Rotacja (orientacja kontrolerów i hełmu) niesie więcej informacji biometrycznej niż pozycja,
        a obie są komplementarne. Pozycja dodatkowo koduje urządzenie (RSA z urządzeniem do 0{,}09),
        czyli cechę niepożądaną w biometrii.
  \item \textbf{Urządzenia:} najwięcej informacji biometrycznej dostarcza Quest (najwyższy Rank-1 w
        9/9 konfiguracjach S2), potem Vive; najmniej -- Cosmos (najniższy Rank-1 w 9/9 konfiguracjach
        S2 i we wszystkich modelach S4; w S4 to 0{,}93--0{,}94 wobec 0{,}97--1{,}00). Większość
        błędów najlepszego systemu pochodzi z Cosmosa.
  \item \textbf{Porównanie klasyfikatorów.} W ustawieniu bazowym (CE): Transformer $\geq$ ResNet1D
        $\gg$ GRU. Transformer najlepiej separuje osoby przy 6 razy mniejszej liczbie parametrów niż
        ResNet1D, a GRU jest najsłabszy i najmocniej koduje urządzenie. W ustawieniu S4 różnice
        prawie znikają, a najlepszy pojedynczy model to InceptionTime. Ensemble różnych architektur
        jest najlepszy (Rank-1 0{,}977, EER 0{,}012), bo modele mylą się na różnych próbach.
  \item \textbf{Metodologia oceny.} Walidacja w obrębie dnia~1 jest nasycona (Rank-1 0{,}97--1{,}00)
        i nie przewiduje wyniku cross-day, dlatego model wybieramy po val loss. Dobór przepisu na
        części dnia~2 zawyża wynik o ok.~1~p.p. (dev 0{,}992 vs test 0{,}977), dlatego raportujemy
        odłożony test. Ranking Rank-$k$ musi być liczony po osobach, a nie po próbkach.
\end{enumerate}

\paragraph{Ograniczenia.} 41 osób i jedno zadanie ruchowe; identyfikacja w zbiorze zamkniętym
(wszystkie osoby są w galerii); dwa dni nagrań, więc nie wiemy, jak wynik zachowuje się po
tygodniach. Przepis S4 wybieraliśmy na 2 seedach i części wariantów, a różnice rzędu
$\pm 1{,}5$~p.p. na dev (do ok.~9 prób z 615) są na granicy szumu. Hipotezy o przyczynach słabości Cosmosa
(jakość śledzenia) nie da się zweryfikować na tym korpusie.

\paragraph{Dalsze kierunki.} Weryfikacja w zbiorze otwartym (osoby spoza treningu w galerii i
próbach), adaptacja domeny między dniami, oddzielne modele lub normalizacja per urządzenie dla
Cosmosa oraz analiza, które fragmenty rzutu (np. przez mapy uwagi Transformera) niosą najwięcej
informacji biometrycznej.
